% =====================================================================
%  R. Nielsen, CHARACTER AND WILL
%  In: For Idea and Reality.  A Journal, edited by R. Nielsen, Bj. Bjørnson and Rud. Schmidt.  Year 1872. First Volume (Bind 6).
%  Copenhagen: Otto B. Wroblewsky, Publisher; Printed by G. S. Wibe.  Bind 6, pp. 489-517.  (Karakter og Villie, in For Ide og Virkelighed.)
%  ENGLISH TRANSLATION, from transcription.tex (Danish), the text of record.
%  GENERATED by ebook-method/for-idee-1870-73/tr_build.py --write from this template and the translators' parts in
%  ebook-method/for-idee-1870-73/tr/ (TRANSLATOR-BRIEF.md; terms: GLOSSARY.md, its key terms approved by Hans
%  2026-10-13/14).  Conventions: TRANSLATION-PLAYBOOK.md.  Page numbers of the print in the margin.
% =====================================================================
\documentclass[12pt,a4paper,openany]{book}
\usepackage[utf8]{inputenc}
\usepackage[T1]{fontenc}
\usepackage{libertinus}
\usepackage{libertinust1math}
\usepackage{textalpha}
\usepackage[english]{babel}
\usepackage[protrusion=true,expansion=true]{microtype}
\usepackage{setspace}
\usepackage[margin=1.3in]{geometry}
\usepackage{fancyhdr}
\usepackage{hyperref}
\hypersetup{pdftitle={Character and Will}, pdfauthor={Rasmus Nielsen}, hidelinks}
\setlength{\headheight}{15pt}
\pagestyle{fancy}
\fancyhf{}
\fancyhead[LE,RO]{\thepage}
\fancyhead[C]{\textit{R. Nielsen: Character and Will (1872)}}
\setstretch{1.3}
\usepackage{marginnote}
\renewcommand{\thefootnote}{\arabic{footnote})}   % the print numbers each page's notes 1), 2)
\newcommand{\opage}[1]{\marginnote{\footnotesize\textit{[#1]}}}
\newcommand{\essayhead}[2]{\begin{center}{\large #1}\ifx&#2&\else\\[0.3em]{\normalsize #2}\fi\\[0.5em]\rule{3em}{0.4pt}\end{center}\medskip\noindent}
\newcommand{\signature}[1]{\par\smallskip\hfill #1\hspace*{2em}\par}

\begin{document}

% the title page of Bind 6 (1872-1-for-idee-og-virkelighed.pdf), in English
\begin{titlepage}
\begin{center}
\vspace*{3em}
{\Huge For Idea and Reality.}\\[1em]
\rule{8em}{0.4pt}\\[1em]
{\large A Journal,}\\[0.8em]
{\small edited}\\[0.4em]
{\footnotesize by}\\[0.6em]
{\large R. Nielsen,}\\[0.4em]
{\large Bj. Bjørnson and Rud. Schmidt.}\\[1em]
\rule{8em}{0.4pt}\\[0.6em]
{\large Year 1872. First Volume.}\\[0.6em]
\rule{8em}{0.4pt}\\[4em]
Copenhagen.\\[0.4em]
Otto B. Wroblewsky, Publisher.\\[1.5em]
{\small Printed by G. S. Wibe.}\\[3em]
{\small Translated from the Danish}
\end{center}
\end{titlepage}

\essayhead{Character and Will.}{}
% --- p. 489 ---
\setcounter{footnote}{0}%
\opage{489}Considered zoologically, man stands on the highest rung
in the series of animals and takes his place of honor ahead of
apes and bats. Whereas older naturalists exerted themselves almost
beyond their strength to make the chasm between
man and the dumb animals deep and wide, the more recent
science has, with particular fondness, pointed out and impressed upon us the
natural kinship between animals and men. To be sure,
it cannot occur to any reasonable student of nature to
deny man an excellent endowment in comparison
with the other animals, but the question is how matters
really stand with this endowment, whether it points
to something original or is merely a difference of degree. That
a man has clearer consciousness and more comprehensive
understanding than a dog, a horse, a fox, an elephant, not
to speak of bats, is beyond all doubt; but
if we now compare a fox, which by everyone's opinion has
what one must call, in the animal sense, a good head,
with an oyster, which is literally headless, then here
again is a considerable distance. And yet Nature maintains
myriads of individuals who, as far as intelligence is concerned,
stand on a lower rung than an oyster. Let us now suppose that
% --- p. 490 ---
\setcounter{footnote}{0}%
\opage{490}all living beings on the globe are originally of
one kind, that the organisms from the beginning have more or less
resembled one another, so that the typical differences
were developed only gradually in the course of the ages;
let us imagine that millions upon millions of
years ago there was a period when the earth's foremost
creature was a fish, a strangely armored, ambiguous
fish; that the breeding of fishes took such remarkable
upswings in the course of millions of years that the ambiguous fish, when it at last
died, left behind a numerous family of fish-lizards, flying lizards,
turtles, iguanodons and crocodiles, which likewise
were ambiguous; that then again, after the lapse of
millions upon billions of years, a period set in when
the old four-legged creatures begot new four-legged creatures, the four-footed ones we
call mammals; that among these four-footed creatures there prevailed a
striving so many-sided that some, which strove upward, became
birds, while others, which strove downward, became moles;
that at the very last there appeared a mother-animal so
doubly pregnant that her young, even before birth,
went in two directions, so that some got no hands
and became guinea pigs, but others got no tail and became
men: then we have, in rough outline, a picture of
the gradual development.

It is by no means our intention in the present context
to go to law with the naturalists about the gradual formation of genera
and species in the history of the earth's development;
we hold to our old principle that the
``heathen, puffed-up sciences'' (!), independent of ``dogmatics,''
must have full right to set themselves their own
tasks and to pursue their own goal by their own rule.
Should the physicists feel the need of a higher rule,
% --- p. 491 ---
\setcounter{footnote}{0}%
\opage{491}should the thirst for knowledge tempt them to take dogmatic
glimpses into the Garden of Eden, in order to enjoy the sight of the
two ``absolutely dissimilar trees'' and to see how Adam looked
before the beasts of prey yet went on the prowl: then
philosophy regrets that it has no objective hints to offer; the
thirsty for knowledge will, in so important a matter of inquiry,
have to seek counsel from bishops and theologians.

Gratifying as it is that knowledge of the
most important truths of natural science is spread and taken up into the general
cultivation, just as lamentable is the superstition of the present age
that the understanding of the inward, psychic life
is to be subjected to a scientific yardstick of natural science, that religion
is to be grounded on zoology, and conscience derived
from physics. It must be impressed again and again that
science itself is innocent of this misunderstanding. True
science keeps strictly to the law of logic that
every conclusion must correspond exactly to its premises.
From this it follows that premises that belong exclusively to natural science
cannot possibly lead to conclusions whose premises
belong to the science of spirit. The science of spirit takes its point of departure
in the inner life of man, his basic consciousness, his essential
self; but man's inner, essential self is just as
little an object for zoology as for geology. It
must be impressed that the modern confusion of spirit, whose
fibrous roots branch out in an uncanny mixture of superficial
considerations and false conclusions, stems from
a faulty dealing with objective knowledge; and so it has, from
the very first, its ground in the old theological deception,
which merely changed its look when one gave up the authority of the Pope
in order to submit oneself to that of science.

The inward world has its facts just
% --- p. 492 ---
\setcounter{footnote}{0}%
\opage{492}as fully as the outward. It is a clear, unassailable
fact that man, among all living creatures
on earth, is the only one that has a judgment about
itself, the only one that can give itself a name and say within
itself: I am a man. In this judgment lies the claim that
man is a self, an I. With this word: I, there
is set an infinite chasm between the cleverest animal and
the stupidest man. To be sure, a clever animal can under
certain circumstances show greater cleverness than a thoughtful
man, and in a natural sense even be said to
stand far above an idiot; but the dazzling quality in this kind of
comparisons is easily removed. When one wants to investigate
the difference between animal and man, it is really
a perverse procedure to set an animal
with healthy organs beside a man with sick ones,
or to regard the certainty the animal has in its instinct
as of one kind with the assurance man must first
fight his way to through reflection and many-sided deliberation.
No, the question is something else: what is the animal
and what is man, when we look at the original
dispositions? The animal is a soul without spirit, the human soul
is essentially spirit; from this point of view the natural boundary is
clear, unambiguous; the animal is sense-consciousness, man's
essence is self-consciousness: this boundary-mark stands firm
so long as the animal is animal, and man man.

But, one says, appealing to geology, the whole
animal kingdom is, after all, in a gradual development, a quiet,
slowly upward-advancing movement; in millions upon millions
of years the animals will found states, the sheep hold
a council of state, the horses study geology, the hens emancipate
themselves, and the cock, crowing, say I to himself. That is a
% --- p. 493 ---
\setcounter{footnote}{0}%
\opage{493}dark saying, which dare not be understood literally. In the first
place, geology has no promise that the living
animal forms will in so distant a future still
exist on earth; on the contrary, it leads us to the conjecture
that the species and genera we know must, when
the great turning-point sets in, all be extinct.
It will then quite certainly not be these animals that reach up to the
human, but other animals, other horses, other hens,
horses with spirit, hens with self-consciousness! In the second
place, the hypothesis of the gradual transition from animals
to men, as regards its plausibility, is as round as
a square. For let us suppose that a time came, if
it cannot happen at once in millions upon millions,
then in billions upon billions of years --- geology has
time on its side, after all --- when the animals really got self-consciousness,
what then? Then the gradual transition is burst asunder, there
has occurred a transformation: the animals, however they otherwise
look, are not animals, but men. Law-giving sheep
and self-conscious hens are nonsense; horses with spirit must
be more than horses; how, pray, should horses be able to
become geologists?

How infinitely deep the difference is between the soul of the animal and
that of man, anyone can convince himself who is willing to
lay any weight on character and will. To be sure, the animal has
character, regrettably not seldom a more consistent
character than a man; but the reason is easy to discover.
The animal's character is a determination of nature, lying
in its species; the animal is from the hand of Nature what it
shall and ought to be. If to live in accordance
with one's nature is virtue, then all animals are virtuous; that is
more than one can say of men. Those who do not
% --- p. 494 ---
\setcounter{footnote}{0}%
\opage{494}rightly catch sight of the actual difference between
man and animal, as far as the positive is concerned, would do well to
begin with the negative. For man can do
what the animal cannot, enervate himself and squander his
life; man can hang himself, shoot himself, ``rip open his own
belly,'' and in general display, as the textbook expresses
it, ``a cruelty toward himself, of which
there are no examples among the wildest animals''; here is a
radical difference. Man has by nature no character,
but certainly dispositions toward character; his peculiar vocation
is that he is, with self-conscious will, to develop
his dispositions and give himself his character: character is
man's real self, and the real self is in the will.

There is a whole circle of questions of life whose treatment
is botched in the common textbooks, because their
possible solution rests exclusively on the formation of character.
Thus the question of the freedom of the will. Those who maintain
that the will, independent of all natural impulses, can
choose, resolve and act as it wishes, proceed blindly.
If it were so that the will could free itself
from all motives, it would have to lie in man's
power to remake his being from the ground up and to make his
mind and his innate dispositions into what they should be.
Experience speaks too emphatically against so wild a doctrine of arbitrariness
for us to waste any time on
refuting it. But to go to the opposite extreme
and to make man's will a nature-bound expression of
an innate disposition of mind, which only upbringing, habit and the example of others
can at most force into a certain
form and guide in a certain direction, is still more dangerous, for in
this way of looking at things there is something deceptive. That the will
% --- p. 495 ---
\setcounter{footnote}{0}%
\opage{495}must have a natural ground, and that this natural ground may
differ in different individuals, is unmistakable.
We see it in children. A child of five
or six years may in certain lights appear to the attentive
observer as a well-equipped little person
who is in the process of giving himself a character; a stubborn
or a soft mind, vehemence, steadiness, boldness, timidity,
openness, reticence, in short, the most diverse
appearances, now of something evil, now again of something
glorious and good, come forth, often in a highly surprising
manner, on the surface of the mind from an obscure ground
and in manifold shadings. Since every spiritual
activity presupposes a natural endowment, and a natural endowment
can often express itself very early, there are
those who think that the will too must rest on a peculiar
natural gift, that the strong mind necessarily must yield
a strong will, the weak mind a weak will, indeed that one
can see already in the individual at a very young age whether
he will ever get a will or is destined to
go through life in the leading-strings of circumstances.
Here we have once more one of those half-truths which, by
being passed off as whole, become untruths. So long as one
one-sidedly regards man from a zoological standpoint,
it is quite true that it is with the
instinctive will as with natural courage. One has
the nature of the lion and is brave, another the nature of the hare and
is cowardly; one got strong drives, great passions and became
a man of will, another got only a drowsy drive,
a languid mind and a lame will: it is, one says, determination by nature
all of it. It is not seldom heard that the man
who is convinced of an evident wrong involuntarily
% --- p. 496 ---
\setcounter{footnote}{0}%
\opage{496}exclaims: ``I will not hear anything at all; it is my nature;
I am once and for all like this and can never become different.''
There is in this sort of defense a peculiar mixture
of defiance, despair and excuse, which bears witness
that the person regards himself zoologically.

The zoological view of man's will stands
in close connection with the statistical one. From statistics too
one derives proofs against the freedom of the will, and the proofs
are so much the more powerful as they can be
built straight upon facts and conducted by sure calculation.
When one has a table of the actions which
over a certain series of years have repeated themselves, in accordance
with the character of the population in a certain country, then one
can predict with certainty what the free will
of this population must will to accomplish in the coming
year. One can predict not merely in general how many
suicides, how many marriages and how many
thefts will occur, but one can go into particulars
and calculate, as regards suicide, how many are
to hang themselves and how many are to drown themselves,
and in which months the free individuals will prefer
hanging to drowning and vice versa; as regards marriages,
how many will marry widows
and how the number of those who marry old
widows will stand to the number of those who marry
young ones; as regards thefts, how many thefts from fields,
how many burglaries and how many pickpocketings
the thievish will will carry out in the course of the year,
and so forth. But if in statistics one can calculate in advance
marriages, suicides, thefts and the like with the same certainty
with which in astronomy one can calculate solar and
% --- p. 497 ---
\setcounter{footnote}{0}%
\opage{497}lunar eclipses, then there seems to lie in this a striking
proof that man's character is on the whole nature-bound
like the animal's, that the so-called human freedom
moves according to laws as necessary as the globes
in the heavens.

To get rid of this whole confusion of geology,
zoology and statistics there is only one way out:
to break off, begin afresh, take the question of the freedom
of the will for what it is: not a question of science,
but a question of character. If the difference between
man and animal is to be recognized as more than a difference of degree,
then it is not enough that man conceives himself as
thinking and falls to musing over the remarkable likeness
that there is between an animal and a man. Man
must begin from the ground up by reflecting on his
original vocation, pull himself together and seek his actual
self in the will. What thinking is we learn by
thinking; but what will is we learn only by willing.
To will is the self's innermost act, the act by which
it raises itself up out of its natural ground, its possibility,
produces itself, makes itself into something real, peculiar,
independent. This act does not
at once aim at getting this or that accomplished in the world,
it is not directed toward the outward, the objective, but
unconditionally toward the inward, the subjective.

There are in the present age two dangerous rocks, on which the individual
may easily founder when he is in earnest about
using his will and becoming a man, the one is a
misunderstood practice, the other a misunderstood theory.
The practical misdirection is that the will, without deeper self-reflection,
turns outward and hastens to get something
% --- p. 498 ---
\setcounter{footnote}{0}%
\opage{498}accomplished; the power of action is sacrificed on the finite in temporal
business and many kinds of civic affairs. Thus
there comes about something that resembles a self, a character,
and yet it is hollow within; for this self does not know
what it itself is, whether it is really something in itself or
merely an organic product of muscles and nerves. It is
uncanny to see how a man can exert
himself to the utmost to show the world how highly he
loves his own self, and when everything then succeeds for him, and
he with his dear self quite imagines himself to be on top,
then he has no self after all. Where did it go? It
was lost to him; it dwindled away in business, it
perished in the outward. But if the practical misdirection is
dangerous, then the theoretical is still far more dangerous.
Theory wants the thinking being not to hasten
with the outward, but to sink deep into itself and reflect
on what will is, before it begins to will.
To will is to think: the pure willing is the pure thinking.
So the thinker sits down and thinks: first the
objective and then the subjective, first necessity and
then freedom, first the zoological and then the psychological;
now thought is in metaphysics, now in ethics, now in history,
now in statistics; and what comes out of the
whole? That the subjective is objective, that freedom is necessity,
that the self is a universal thinking. Theory
is consistent; it ends as it begins: it begins
with thinking and ends with thinking. The objective,
universal thinking is everything, has everything, takes everything, drowns
everything, the thinker, the self, the will, in eternal necessity;
then the will is thought, then the thinker can begin to will
--- to think.

% --- p. 499 ---
\setcounter{footnote}{0}%
\opage{499}Only when the question of freedom is transformed into a question of character
is everything clear. Man is not
man in the same sense in which the animal is animal; man
is, by his will, to make himself into what he essentially
is, into a real man. As regards the real man,
it is not the face, the two hands and
the upright gait that are decisive; no, it is the will:
such a will, such a man, so much will,
so much man, neither more nor less. From
this point of view it is then also easy to come to terms
with zoology. The disposition of the will makes a difference, a difference valid for time
and eternity, between man and animal,
but no difference between man and man, for
every man has the gift. Powers of mind and passions
are highly different in different men; but
willing, or rather: being \emph{able} to will, is common
to all men. Self-consciousness, self-determination
and will are in their disposition one and the same. Here it is not asked
how great the natural force of the will is, for the instinctive,
nature-bound will stands undeniably in proportion to
individual endowment; but here it is asked about the inner
essence of self-determination, about the will's exertion to
work itself up to real will. Let a man be
ever so frail and have fallen into ever so many temptations;
if such a man could say to himself with a good conscience:
I know that I have honestly summoned the whole strength of my
will and resisted to the very utmost, then
this man, in the midst of his temporal abasement, would stand
as an exalted self with a will more heroic than
Alexander's, Caesar's or Napoleon's. The fact is that, when it
is a matter of willing to do what the self recognizes as
% --- p. 500 ---
\setcounter{footnote}{0}%
\opage{500}noble and good, or of willing to resist what it recognizes
as base and bad, then the human will never goes
to the utmost. This weakness, an incomprehensible self-contradiction,
is common to the hero and the most insignificant man.
The more developed the consciousness of will is in a
man, the more deeply he feels that he has made himself
into a more or less botched character, who bears
a responsibility; the less there is in an individual of a real
man, the more easily he is satisfied with
excuses; the more deeply developed the man is, the
more he will have to reproach himself with. Self-reproach,
this strict boundary-mark between animal and man, is a
telling testimony to the originality of freedom, to the will's
inner infinity. From this point of view it is easy
to come to terms with statistics. Statistics does not deal
at all with what men can do, but with what
they are wont to do; it does not single out the individuals
and assert that these individuals, each for himself, inevitably
must will this or that, on the contrary! about the individuals
as individuals statistics knows nothing. It takes individuals
by the multitude, regards them as atoms in the mass,
and so arrives, quite correctly, at an average result
which clearly shows that individuals en masse, civilized as
well as uncivilized, are as yet not whole, but half men.

That the question of the freedom of the will is most
closely connected with the question of religion and faith
no one will surely deny. If now the question of freedom, taken originally,
is not a question of science, but a question of character,
then the same must surely hold of the question of faith.
``Show me thy faith by thy works,''
this proposition, which, although its truth is both great
% --- p. 501 ---
\setcounter{footnote}{0}%
\opage{501}and clear, has aroused so many disputes, misled so many
souls and confused so many brains, may at last
be freed from all misunderstanding and set in its
proper light, when it reads: ``\emph{Show me thy faith by thy
character}!'' Character is something more than a sum of
works, and yet works belong to the mirroring
of the real character; works are like lineaments
in the physiognomy of character, but character itself
is inwardness, is spirit and essence. If it has hollowed out Catholicism
and struck Pietism with spiritual unfruitfulness
to want to hang everything on works, then Lutheran
orthodoxy has surely likewise had to suffer for wanting to keep
works outside the essence of faith. Both wrong paths have
their point of departure in the theological disputes about the law, in
which a scholastic acuteness has been tempted
to separate what absolutely must be united in the religious
character. It is perfectly true, what Luther
has taught, that faith is essence, that this essence is life, that
life in inwardness is a power of God unto salvation; but
when the powers of inwardness stir, then works
come forth, and then we have character. If the demand for character is left out,
consciousness generally divides
into two halves; the one half belongs to eternity, the
other to temporality, and then one can again go different
ways. If it does not depend on works at all, but
only on faith, then it lies near at hand to conceive faith
as a matter by itself; a disposition of mind that accepts the pure
doctrine, lets itself be edified by the ceremonies
ordained in a Sunday church, and beside this has its works
in the temporal, attends to its office, its trade,
lives with all civic rectitude in the civic
% --- p. 502 ---
\setcounter{footnote}{0}%
\opage{502}outside the eternal. It has rightly been maintained that
``the Christians'' nowadays are not to be distinguished from other
decent people. What difference is there, pray, in
the cultivated world between Jews and Christians? The Jews
go to the synagogue on Saturday, the Christians go to church
on Sunday; that is the whole difference. At the pinnacle
of the present European civilization there are ---
thanks to classical intelligence --- neither Christians nor
Jews nor heathens: we are all neutral men of culture.
But when intelligence dares not give up religion, although
the faith, separated from works, is too inwardly feeble to
give itself real character in a real world, then the
consequence is a spiritual deadness, which is all the more
destructive to the inward man as it not
seldom takes pleasure in the conceit that one nevertheless
really has religion, when one vigorously holds to the state church
and cherishes a vague respect for the dogmas. How far
such a deadness can go in an old culture,
and how easily it lets itself be combined with a conceited regard
for something holy, ideal and exalted, of this the following
example is a telling proof.

The missionary Huc relates\footnote[1]{\textit{L'empire chinois} 1, 173; cf.\ Oscar Ræder: ``Historical Laws''. Copenhagen 1871, pp.\ 20---21.}, how exceedingly difficult
it is to make proselytes for Christianity in
China, since indifferentism, the religious characterlessness,
is so great that nothing can overcome it. It is far
from the case that the missionary is received only with anger; on the contrary!
The Chinese listens to the stranger with pleasure; it
interests him and awakens his thirst for knowledge. He
% --- p. 503 ---
\setcounter{footnote}{0}%
\opage{503}accepts, approves, all that is said to him, and raises not the
slightest difficulty or objection; he finds that what is said is true,
beautiful and excellent, indeed he himself appears as a preacher, and now one hears him speak
of plundering the images of the gods in favor of Christianity.
He laments the blindness of the people who
attach themselves to the goods of this world, and he would, if
it were asked of him, deliver a splendid speech on the happiness of
knowing the true God, of serving him, and in this way
winning eternal life. Hearing him, one would
think that he was already a Christian, but he has not come
a step further. And yet one must not think that
his words are not, to a certain degree, sincere; he
believes what he says, he likes to talk about it, but always
as of a thing that does not exist for him and does not
concern him. The Chinese admirer of the Christian
faith then sometimes ends by declaring to the missionary, who
begs him to hasten his conversion, that, though he
is not blind to the beauty and sublimity of the Christian religion,
he nevertheless does not see why one should increase the
afflictions and cares of life more than necessary by occupying
oneself with another life than the present one and
striving beyond one's capacity. The body that we see and
touch costs us every day the greatest attention and care;
should we then on top of that occupy ourselves with a soul that
we do not see, and a life that we do not know? --- Poor
Chinese, he has faith, but lacks merely a couple of works.
Had he been born in the maternal womb of a European Protestant
state church, then the works would have
come as if of themselves; for he has the gifts, after all.
Imagine that such a man had studied
% --- p. 504 ---
\setcounter{footnote}{0}%
\opage{504}theology, and learned to preach in the Pastoral Seminary: he would have become something great; he would surely have become more than a dean,
he would have become a bishop and a witness to the truth.

By letting the works go to the temporal and
faith to the eternal one can also enter upon
another wrong path, closest to the Pietist. The feeling
presses itself upon one that here in the temporal one really must
do something for the eternal; but what is one to do?
The temporal is temporal and never becomes anything else. So
one must begin by renouncing the temporal. How much
is one now to renounce? Not everything temporal; for if
one literally renounces everything temporal, one dies
at once; one takes one's own life. So one divides
temporal things and works into two classes. The first
class comprises the necessary and permissible works,
such as: eating, drinking, sleeping, founding a family, pursuing a civic occupation, etc.; the temporal contained in these is
accordingly not to be renounced. To the second class belong not
only such vices and bad habits as the police punish,
but also actions that moral society otherwise
regards as innocent and permissible, e.g. dancing, playing
cards and going to the theater, indeed, what is more: loving
one's people and defending one's fatherland, grieving at its misfortunes,
rejoicing in its honor, honoring its memories and working
together for its future, are reckoned among the follies
by which the Devil in particular weaves his snares, when he goes about
craftily and catches souls in the temporal. But who now
does not see that a register-like enumeration of what is to be
renounced and not renounced is better suited to hatching
screwed-up individuals than to cultivating characters. We
honor the Pietist earnestness when we warn against its
% --- p. 505 ---
\setcounter{footnote}{0}%
\opage{505}mistakes. It is, with regard to the great, perhaps terrible,
crises of the spirit that the future holds in prospect, a true
good fortune for our fatherland that it possesses in its literature a
writer who, from a Pietist standpoint, has in the
most comprehensive and thorough manner dissolved
the whole of Pietism; that is S. Kierkegaard.

To break with the world, renounce the world, die away from the world,
this is and remains, as long as there is Christianity,
a Christian basic demand, an unrejectable demand of character,
inasmuch as by the world is then understood what we all know,
the sensual, selfish, temporal, vain being. The demand
is negative, an unmistakably genuine demand of the Law.
Wherein, then, does the Pietistic consist? Not in this, that
``the ideals are proclaimed,'' that the demand is made to hold
in its whole extent; for that is precisely the Christian,
and anyone who, under the mask of mildness, on behalf of softness,
permits himself to keep silent about, tone down or veil
the holy demand is certainly not a Christian
character, still less a witness to the truth, but a
traitor who kisses the Crucified. No, the Pietistic
consists in something else; it consists in this, that
the negative demand of the Law is put in the place of the Gospel,
set up by itself, as if it itself contained the conditions
for its fulfillment. This is what is unevangelical.
Now it lay not only in S. Kierkegaard's exclusively
nervous individuality --- either \emph{obedience} in faith or a
place in the madhouse --- and in his quite peculiar
development of character, but also in the way in which he
conceived his mission, his struggle against what he regarded
as an Epicurean, enervated age, bewildered in worldliness, stupefied by
``priestly twaddle,'' that he put the demand as
% --- p. 506 ---
\setcounter{footnote}{0}%
\opage{506}he did. Thereby it came about that his presentation of
Christianity, although it indisputably contains all the basic determinations of the Gospel,
took on a peculiar sheen, a touch of
Pietism.

But the Pietists will hardly rejoice when they consider
how their renunciation of the world is here carried through.
To set religious earnestness in abstaining from
plays, cards and dance, whining over the ungodliness of the world
and for the rest being like the others --- a bit of a thief in
one's trade, is not earnestness; when such a renunciation wishes
to pass for more than what in many cases it
may well be, an advisable caution against falling into temptation,
wishes to count as full earnestness, holy, deep and godly
earnestness, then this sort of godly earnestness, as anyone
can see, is simply ``handed over to comic treatment'' by S. K.
No, the demand is that all, unconditionally all, ---
the means of livelihood too --- shall be renounced. The demand is that
the interests of temporality shall all be burned up in the fire of the spirit,
and the selfish consumed by the flame. The unconditionality
in the demand, that all, all temporal things without exception,
shall be renounced, is the equal, and this equality is
again what is exalted in the earnestness.

But let us now see how the exalted earnestness
looks when the ideals are proclaimed and a human
will is to try to set the eternal into character.
``Wherever the Christian is,'' says S. Kierkegaard in
``Works of Love'' (p. 63), ``there is also self-denial,
which is the essential form of Christianity\dots{} But
erotic love and friendship are precisely the highest of self-feeling, are
the I intoxicated in another I. The more firmly the two I's
join together to become one I, the more
% --- p. 507 ---
\setcounter{footnote}{0}%
\opage{507}this united self shuts itself out selfishly from all others. In erotic love
and friendship at their highest the two really become one
self, one I. This is explicable only because in preferential love
there is contained a determination of nature (drive --- inclination)
and self-love, which can selfishly unite two
in a new selfish self. The love of the spirit, on the contrary, takes away
from my self all determination of nature and all self-love,
therefore love of the neighbor cannot make me
one with the neighbor in a united self\dots{} Are you, my hearers
(p. 67) perhaps, as the saying goes, a cultivated person, well then, I
am also cultivated; but if you believe that by means of
`cultivation' you come nearer to this highest, then
you are gravely mistaken. And here precisely lies the error; for
cultivation we all wish, and cultivation has the
highest constantly in its mouth, indeed, no bird that has learned but one
single word calls out more incessantly this single word,
and the crow no more incessantly its own name than
cultivation always calls out `the highest.' But the Christian
is by no means the `highest' of cultivation, and the Christian
cultivates precisely through the repulsion of offense\dots{}. You
\emph{shall} love, you shall love your neighbor as yourself! (p.
70---71). Do not refrain from believing, because the commandment almost offends
you, because the speech does not sound ingratiating like the poet's,
who with his songs ingratiates himself in your happiness...
No, love the beloved faithfully and tenderly, but let love
of the neighbor be the sanctifying thing in the covenant
of your union with God; love your friend sincerely and devotedly, but
let love of the neighbor be what you learn from each other
in the confidence of friendship with God!.. 
As the glad tidings of Christianity
are contained in the doctrine of man's
kinship with God, so its task is
% --- p. 508 ---
\setcounter{footnote}{0}%
\opage{508}man's likeness to God\dots{}. Insofar as you love the
beloved, you do not resemble God, for with God there is no preferential love\dots{}.
Insofar as you love your friend, you do not
resemble God, for with God there is no distinction. But when
you love the neighbor, then you resemble God.''

Anyone who earnestly wishes to take to heart how
the opposition between self-love, the unchristian,
and love of the spirit, the Christian, is here presented,
must necessarily become aware that the presentation
contains two explanations which explain one another in such a way
that light with light gives darkness.

First explanation. The Christian has no beloved,
no wife, no friend; he has only the neighbor! If anyone
points out to him that he must necessarily
stand in a different relation of love to his beloved, if
he is betrothed, than to other girls, to his wife, if
he is married, than to other wives; it does not help. If
he really is in the character of what is here set forth as
the Christian, he will love all girls, all men and
all wives in the whole world equally, in the same spirit, with
the same inwardness and in the same way, for he loves
in them all --- the neighbor. If anyone reminds him that
besides the universal love for all human beings
he must necessarily cherish a quite special love
for certain human beings, that in his beloved he is indeed to love
his neighbor, but yet also conversely in \emph{this} neighbor his
beloved, in \emph{this} neighbor his friend: it is in vain. The
eye with which he should see the beloved as the beloved,
the friend as friend, has been put out; it was a heathen
eye, after all; so they are both, the beloved and the friend,
gone for him: in their place stands the neighbor.
% --- p. 509 ---
\setcounter{footnote}{0}%
\opage{509}Even if he could discover them, how should he
be able to love them. Love of the beloved is preferential love,
love of the friend is preferential love; preferential love
is self-love, self-love is selfish, and
the selfish ungodly.

Second explanation. Although the Christian in all
human beings loves only the neighbor, he nevertheless has permission
to have a friend and a beloved, whom he may love in particular.
``Love the beloved faithfully and tenderly, but''....
``love your friend sincerely and devotedly, but''\dots{}. The meaning
cannot possibly be that the beloved and the friend are merely
to be loved as the neighbor; in that case the repeated
But would be without all significance. The permission, if one
will, the admonition must be explained thus: If you, the
Christian, only do not forget to love every human being without
distinction, your friend, your beloved too, as your neighbor, then add
just so little preferential love: love --- with preferential love
--- your beloved faithfully and tenderly; love your friend with
sincere and devoted preferential love! But what was
preferential love according to the first explanation? Selfishness,
sin, ungodliness! ``Insofar as you love the beloved,
you do not resemble God, for with God there is no preferential love.
Insofar as you love your friend, you do not resemble God, for
with God there is no distinction.'' And what is preferential love
according to the second explanation? A little indulgence, which the author
offers us, not for money, that would be a sin to say,
but gratis and from a good heart. But is it not
--- from the standpoint of the Law --- as if the author said to his
earnest reader in plain words: Now you know what sin
is, that it is love of distinction, preferential love. But
when you know it, acknowledge it, feel it and bow
% --- p. 510 ---
\setcounter{footnote}{0}%
\opage{510}deeply before the ideal; then sin, only --- ``faithfully and tenderly'' with
the beloved, ``sincerely and devotedly'' with your friend.

That it cannot occur to any sober-minded person
to want to foist on S. Kierkegaard so insane a distortion
of the Holy hardly needs to be remarked;
nor does the contradiction lie in the demand of the ideal
as ideal. No, the basic contradiction, the icy, cutting
contradiction, arises from this, that the Law immediately
is to demand of the believer's will that it be set into character
as law. This is what, for all godly talk of
regeneration, sealing and sanctification, is the Pietistic.
To unfurl the banner of the ideal with the inscription of the Law is to
dissolve Pietism.

We now turn to the Grundtvigian view, to
see how matters stand here with character and
will. Grundtvig does not begin with the Law, but with
the Gospel; he does not speak of faith and works, but
of the power of the spirit in ``the living word,'' of the renewal of life
in faith; he does not set the demand of the spirit in hostile
opposition to man's determination of nature, but reconciles
both and finds the source of unity in the heart. What
follows from this? The great commandments of the Law: You shall love
God above all things and your neighbor as yourself, are
hereby neither overlooked nor pushed back, but presented
in the light of the Gospel, so that the commanding: You \emph{shall},
is transformed into a promise: It shall be given you to \emph{be able}.
On which side is there now, when what is at stake, mark well, is
not merely the demand but the fulfillment, the
deepest earnestness? Obviously on the Gospel's. The Law says:
You \emph{shall}, you \emph{shall}! and yet knows very well that man
cannot do what he thus shall. But to demand
% --- p. 511 ---
\setcounter{footnote}{0}%
\opage{511}of a being that it absolutely shall do what it demonstrably
cannot is one-sided earnestness. The Gospel says: You shall, for you
can; it is I who give you strength. Precisely the fact that
nothing other or more is demanded of man than what
man essentially can do raises infinitely the earnestness of the
demand. Taken originally, the Gospel does not begin with
a demand, but with an offer. It does not demand of
man that he shall on his own fulfill this
or that law, but it offers a gift, the gift of the spirit
and of love, whereby it becomes possible for man
to fulfill the law of love, still more clearly:
whereby it becomes possible for the divine love
itself to perfect its own law in man. The Gospel
does not address itself with a commandment of duty, a prescription, to
man's will, does not begin by saying: You are obligated
to believe in the Crucified, obligated to
let yourself be baptized, etc. No, the Gospel goes the way of inwardness,
begins with a promise and speaks to
man's heart, for the whole man is in the heart. Does
the Gospel then perhaps make faith a dim feeling, without
spirit and clear thoughts, without character and will? By no means!
But the Gospel does not know, will not know any
final parceling-out of our inner being; on the contrary, it will
heal what is torn asunder, make whole what is divided, and abolish
the piecework: feeling and thought, spirit and will must
in truth be one if man is to attain
uprightness and acquire character. The Gospel does not know, will
not know any inner discord between spirit and nature;
therefore it seeks at once, without circumlocution, the heart, for
only in the heart is spirit nature, and nature spirit. A
nature-hostile spirit and a spirit-forsaken nature are
% --- p. 512 ---
\setcounter{footnote}{0}%
\opage{512}uncanny distortions of life and truth. Will the Gospel then
perhaps lure human beings to itself and win them by
flattering the natural heart? By no means! The Gospel
knows the human heart from the ground up and knows
full well how matters stand in there; but it knows
also that the way of the heart is the only possible
way of love, that there is no other. If it were only a question of
civic probity, of the fulfillment of duties in obedience,
the spirit could easily use force, strike terror
and compel obedience. But when obedience is to go
hand in hand with freedom and have its ground in love,
the spirit cannot use force, for love does not let itself be
compelled. Love can drive out fear, but no
fear in heaven or on earth will be able to
drive love in. The Law can threaten, compel, thunder,
command: You shall, you shall! to all eternity; its threat is
cold and barren and cannot beget love; but the Gospel
begets love --- in man's heart. If
in a human nature there is nothing but selfish
craving, sensual desire and the lust of the flesh, then
Christ never had a human nature, and his
Gospel is an impossibility. But when the Gospel goes
to the heart and makes, in a divine sense, the natural
spiritual and the spiritual natural, then it is earnestness,
holy earnestness; for it is not man's carnal
heart that rules over the Gospel, but the Gospel's spirit
that rules in the heart.

On this foundation ``the ideals can be proclaimed,''
characters built up, the will strengthened as the organ of the spirit
and ``Works of Love'' practiced. But it will then
also be seen that the whole concept of love, conceived from the
% --- p. 513 ---
\setcounter{footnote}{0}%
\opage{513}human side, must undergo a transformation. Self-love
and preferential love are radical determinations,
inseparable from all that human love which
from any side, in any sense, in any way can be said
to have a natural origin. If we pull up the ramifications of self-love
and of preferential love by the roots, then the Gospel departs and leaves behind a nature-hostile
spirit with a beginning devoid of nature. Self-love
and preferential love are ambiguous words; the ambiguity shows
that they can belong to both the lowest and the highest.
Take preferential love away, and what becomes of a mother's love?
Take self-love away, and what becomes
of love of the neighbor? ``Love your neighbor as
yourself!'' But to love oneself is indeed self-love. In the selfish
sense self-love is self-seeking love, and self-seeking love
incompatible with love of the spirit. In order to master
egoism and purge self-seeking love, the spirit of the Gospel
enters into a natural alliance with the deeper self-feeling,
with the self's need for eternal life, that is, with the
true self-love; then the demand of the spirit becomes at bottom
a natural impulse, which awakens the will through concern for oneself,
and the will awakened through concern for oneself becomes the
organ of the spirit. That the instinct of self-preservation is a strong
instinct holds true not only in the sensual but also in the
spiritual sense. Once it has become evident in earnest to a human being
that self-seeking love is his worst
enemy, then he will hate self-seeking love, not because he
hates himself, no, because he loves himself. That the spirit's will
hates self-seeking love and hates it thoroughly, precisely
because it hates out of self-love, means, in other words,
that the will with all its strength combats the selfish, does not
% --- p. 514 ---
\setcounter{footnote}{0}%
\opage{514}stop at that, does not set itself at rest until the fury of self-obsession
is quenched and the deadly enemy strangled. It is a work of the
spirit, and yet it proceeds naturally: in true self-love
the natural is spiritual, the spiritual natural.
With preferential love the matter stands in the same
way. All preferential love is natural and must, insofar as
it keeps to sensual selfishness, necessarily be rejected
by the spirit. But it is only the selfishly narrowing,
the partiality, the bewitchment, the infatuation in preferential love,
not preferential love itself, that the spirit hates. Love
of others is impossible for human beings without preferential love.
Real love must, like everything that is
to become real in the world, necessarily be limited.
An unlimited love, a perfectly equal, universal
love of all human beings in general
is indeed an expansion of the mind in imagination, but as yet no
real love, for it is not love of any
real human being. If love is to become real,
it must limit itself; but in limiting itself and
choosing its object, it becomes preferential love. It goes without saying
that false preferential love must be related
to true preferential love as self-seeking love
is related to true self-love. The husband
is, according to Paul's exhortation, to love his wife as
he loves himself. If the husband's self-love now
is nothing other than self-seeking love, then his love
for his wife, however inward, tender and enthusiastic it
may otherwise look, can likewise be nothing other than a
self-seeking love converted into preferential love, a preferential love
that selfishly excludes love of the neighbor. If,
on the other hand, his self-love is in truth a love of the spirit
% --- p. 515 ---
\setcounter{footnote}{0}%
\opage{515}purified of egoism, then his love for his wife
will also from the very beginning be love of the
neighbor. Let him only begin by having this neighbor
as ``the beloved'' and be as blissfully in love as you like;
provided only that his will is in the character of faith, the spirit of the Gospel
will vouch that the fire that burns in his
soul shall unfailingly purify the sensual being; let
the first immediate natural impulses vanish with the
years, love will nevertheless not cool, but inwardness
will grow with rejuvenating strength and give the character
an expression which, although it is not in the bloom
of youthfulness, is nevertheless for that very reason just as rich, indeed richer,
because it is spiritually transfigured and yet withal not less
natural. If ever a man has preserved and nourished
such a love for his wife, then there is no need here of any
heralds of poetry to trumpet forth his erotic love: religion
acknowledges him, that his life has been in the
service of the spirit; the Gospel gives him the testimony
that he has loved his neighbor as himself. If anyone
objects that a limitation such as the one here indicated
stands far too low beneath the demand of the ideal, that it is not
enough to love a single human being, least of all one's wife,
but that one is to love strangers too, many
human beings, indeed all human beings: then the answer must be that
in this very matter it does not depend on a plurality.
The neighbor has, remarkably enough, no plural. We have
not many neighbors, but only one neighbor, each time in
each real personal relation the neighbor, \emph{this} our neighbor.
The neighbor is not only One in All, but also All in
One; he who in truth loves a single human being as his
neighbor loves from the heart all human beings, every human being
% --- p. 516 ---
\setcounter{footnote}{0}%
\opage{516}without exception, for if in his heart he makes an exception,
then his love, toward the many just as well
as toward the single one, is an egoistic deception. He may for
that matter well love ten human beings; if he
makes an exception against the eleventh, then the whole is
an illusion. Equality aims at possibility, for it
concerns the disposition; but limitation aims at reality,
for it concerns the act. He who has
learned to love his neighbor as himself has much in the world
to love. He loves not only his nearest, his
family, his friends, but also his people, his fatherland.
To want to love one's neighbor without loving one's family is
masked self-seeking love; to want to love one's family and
one's friends without loving one's people and one's fatherland is
false preferential love. To love one's people and one's fatherland
with true preferential love is infinitely different from
being infatuated with one's own tribe and deifying such national
virtues as in all their naturalness are still always paired
with corresponding bad habits and vices; to love one's people
with preferential love cannot, as truly as preferential love is in the
character of spirit, possibly lead to hating or
despising or misjudging the good and excellent that
is to be found among other peoples, since an inward, strong and living
preferential love for one's own people is on the contrary the only
possible way to love the human in real human beings,
the only possibility for love to become again
flesh and blood, for health, strength
and marrow to come into the otherwise so pale, universal love of humanity.

\emph{Show me your faith by your character}! The temporal
is to be renounced, and yet it is to be used vigorously; the
% --- p. 517 ---
\setcounter{footnote}{0}%
\opage{517}Christian is to ``die away from the world,'' and yet he, precisely
he, is to live in the world, work in the world, labor to carry
human life forward to richer, truer, fuller development
in the world. ``The ideals are to be proclaimed,'' the lofty models
of suffering, sacrifice, struggle of soul and renunciation
are to be set forth in that light of glory which shines with eternal
splendor, and the disciplining commandment of the Law: You shall, you shall!
is to sound constantly with all its severity; and yet we
human beings, without fear of the Law, in the spirit of freedom, of self-understanding,
of the Gospel, of love, are to begin ---
rightly humanly with the divine, rightly naturally
with the spiritual, are to work and act in all confidence, with joy
and hope in the imperishable, in the midst of
all this perishability, this temporality, this
worldliness: these are, to a superficial view, contradictory
demands. Show me your faith by your character! for these
demands are tasks, urgent tasks, not for
objective knowledge and reasoning criticism, but for character
and will.
\signature{R. \emph{Nielsen}.}

\end{document}
